\documentclass[12pt]{article}
\usepackage[utf8]{inputenc}
\usepackage[T1]{fontenc}
\usepackage{hyperref}
\title{Venetian Language Resources for K-12 Curriculum Development}
\author{Compiled by Claude Sonnet 5 (claude-sonnet-5) for Andr\'as Kornai}
\date{}
\begin{document}
\maketitle

\section{Introduction}
Venetian (ISO 639-3 vec) is a Romance language spoken mainly in Veneto, with speakers in Friuli-Venezia Giulia, Trentino, Istria and parts of Slovenia, Croatia and Montenegro, and in emigrant communities in Brazil and Mexico. English Wikipedia gives 3.9 million native speakers (2002). This search was limited and did not locate current grammars, school textbooks or the Venetian-language Wikipedia, so the list below is short.

\section{Primary Educational Resources}
No textbook or primer was located. Search found news about a 2022 bill by Lega deputies (first signatory Massimo Bitonci) to include Veneto among the minority languages protected by law 482/1999 and to teach it in school, especially in nursery school. The outcome was not found and should be checked.
\begin{itemize}
\item Orizzonte Scuola report on the proposal. \url{https://www.orizzontescuola.it/lezioni-di-dialetto-veneto-nelle-scuole-la-proposta-della-lega-i-bimbi-lo-imparino-fin-dallinfanzia/}
\item Sky TG24 report. \url{https://tg24.sky.it/politica/2022/11/17/scuola-lega-dialetto-veneto}
\item Money.it report. \url{https://www.money.it/insegnamento-dialetto-veneto-scuola-tramissioni-tv-proposta-lega}
\end{itemize}
A trainer's CV (a partisan, unverified source) mentions classroom projects such as ``Percorsi Linguistici Veneti'' in lower secondary schools. \url{https://www.academiabonacreansa.eu/wp-content/uploads/2024/03/CV-completo-Alessandro-Mocellin-marzo-2024.pdf}

\section{Linguistic and Reference Materials}
\begin{itemize}
\item Wikipedia, ``Venetian language''. \url{https://en.wikipedia.org/wiki/Venetian_language}
\item Boerio, \emph{Dizionario del dialetto veneziano} (1829; later editions 1856, 1867), described on History Walks Venice. \url{https://historywalksvenice.com/bibliography/boerio-1829/}
\item Drusi and Vescovo, article on the history of Venetian as a written language, with discussion of Manin's grammar and Boerio's dictionary (\emph{Quaderns d'Itali\`a}). \url{https://raco.cat/index.php/QuadernsItalia/article/download/26274/26108/26198}
\end{itemize}

\section{Religious and Community Resources}
None found in search results.

\section{Digital and Audio Resources}
None found in search results.

\section{Curriculum Development Context}
The language is stable and widely spoken, but it is not known to be taught in schools, and it has no legal minority status at national level as far as the sources found show. Further work should look for modern orthographic standards, current dictionaries and grammars, and the Venetian Wikipedia, none of which this search confirmed.

\end{document}